\section{Modelado y propiedades ópticas}
\label{section:optical_properties}

Para modelar la morfología bicóncava característica de los eritrocitos, se han propuesto diversos enfoques analíticos, entre los que destacan los modelos de Skalak \cite{skalakStrainEnergyFunction1973}, Lu \cite{luSimulationsLightScattering2005}, Fung \cite{fungHighresolutionDataGeometry1981},
 y los óvalos de Cassini \cite{ereminaDifferentShapeModels2006}. En este trabajo se emplea el modelo basado en óvalos de Cassini, debido a su simplicidad analítica y su capacidad para reproducir fielmente la morfología promedio de un eritrocito sano, optimizando el costo computacional en simulaciones de esparcimiento \cite{ereminaDifferentShapeModels2006}. Los óvalos de Cassini son curvas planas definidas como el lugar geométrico de los puntos cuyo producto de distancias a dos focos fijos es constante; esto los distingue de las elipses, donde se conserva la suma de dichas distancias \cite{wriedtLightScatteringSingle2006}. Al rotar un óvalo de Cassini alrededor de su eje de simetría, se genera un sólido de revolución con la geometría de un disco bicóncavo. En coordenadas cartesianas, asumiendo simetría respecto al eje $z$, la superficie se describe mediante la expresión \cite{ereminaDifferentShapeModels2006}
 %
 \begin{tcolorbox}[ams align]
 	z(x) = \pm c \sqrt{-a^2 - x^2 \pm \sqrt{4x^2a^2 + b^4}},
 	\label{eq:cassini}
 \end{tcolorbox}
 %
\noindent donde los signos $\pm$ externos definen la simetría respecto al plano XY, mientras que los internos delimitan las ramas del óvalo. Los parámetros geométricos se interpretan de la siguiente manera: $a$ representa la semidistancia focal, $b$ es la raíz cuadrada del producto constante de las distancias a los focos, y $c$ es un factor de escala que ajusta el espesor máximo sin cambiar su radio \cite{wriedtLightScatteringSingle2006}. La relación entre estos parámetros delimita la forma del eritrocito: un incremento en $a$ acentúa la depresión central (bicóncavidad), mientras que valores menores de $a$ derivan en geometrías convexas o incluso esferoidales. Por su parte, $b$ escala las dimensiones globales de la curva, controlando así el volumen total \cite{ereminaDifferentShapeModels2006}. En la Fig. \ref{fig:Cassini_ovals} se muestra un modelo con parámetros $a=2.6~\mu$m, $b=2.7~\mu$m y $c=0.813~\mu$m, los cuales corresponden a un eritrocito sano con un volumen de $80.62$ $\mu$m$^3$. La Fig. \ref{subfig:bidimesional} muestra la sección transversal en donde se identifican en color rojo los dos focos ($F_1,F_2$) del óvalo separados uno del otro a una distancia focal de $2a$. Además, se muestra el valor de $c$ sobre el eje $z$ y se muestran en color gris las distancias entre los puntos focales y un punto arbitrario $P$, cuyo producto determina el valor de $b$. La Fig. \ref{subfig:front} presenta el corte del sólido de revolución resultante.
%
\begin{figure}[t]
	\centering
	\sidesubfloat[First image]{\hspace{-0.5cm}{
		\includegraphics[width=0.53\textwidth]{2-Eritrocitos/Cassini.pdf}\label{subfig:bidimesional}}}\hspace{0.1cm}
	\sidesubfloat[Second image]{\hspace{-0.2cm}{	\includegraphics[width=0.35\textwidth]{2-Eritrocitos/axial_ery.pdf}\label{subfig:front}}}
	\caption{\textbf{a)} Óvalo de Cassini con $a=2.6\,\mu\text{m}, b=2.7\,\mu\text{m}$ y $c=0.813\,\mu\text{m}$ correspondiente a un corte transversal de \textbf{b)} el sólido de revolución que modela a un eritrocito sano con volumen de $80.62\,\mu\text{m}^3$. Los focos $(F_1,F_2)$ se identifican con color azul y se encuentran separados a una distancia de $2a$. Se muestra en color rojo el valor de $c$ para el caso particular sobre el eje $z$. Además, se observan en color gris las distancias entre los puntos focales y un punto arbitrario $P$. El producto de distancias de los focos $F_1$ y $F_2$ a $P$ determinan el parámetro $b$. }	\label{fig:Cassini_ovals}
\end{figure}
%

Para estudiar los procesos de esparcimiento y absorción de la luz en eritrocitos es necesario considerar tanto la descripción morfológica proporcionada por los óvalos de Cassini, que define el volumen ocupado por la partícula, como su respuesta electromagnética, determinada por el contraste entre las propiedades ópticas del eritrocito y las del plasma \cite{bosschaartLiteratureReviewNovel2014}. En el caso de los eritrocitos, su función dieléctrica está modulada la hemoglobina, ya su vez por variaciones de la CHCM, lo que modifica la respuesta óptica de la célula \cite{bosschaartLiteratureReviewNovel2014}. Experimentalmente, los valores del índice de refracción de la hemoglobina se han obtenido a partir de mediciones de transmitancia y reflectancia mediante el uso de esferas integradoras y técnicas espectroscópicas \cite{friebelDeterminationComplexRefractive2005, meinkeOpticalPropertiesPlatelets2007}. No obstante, la mayoría de los datos disponibles en la literatura para distintas concentraciones de hemoglobina se encuentran restringidos a ventanas espectrales finitas.

Los datos experimentales de la parte real e imaginaria del índice de refracción de la oxihemoglobina empleados en este trabajo corresponden a tres valores de la CHCM\footnote{La CHCM cuantifica la concentración total de hemoglobina intracelular y no distingue entre oxihemoglobina y desoxihemoglobina. En consecuencia, el uso de propiedades ópticas de la oxihemoglobina es consistente con la definición de CHCM y con su empleo en modelos ópticos de sangre bajo condiciones fisiológicas promedio \cite{bosschaartLiteratureReviewNovel2014}.}: 28.7~g/dL~\cite{friebelModelFunctionCalculate2006}, 30.6~g/dL \cite{friebelDeterminationComplexRefractive2005} y 34.0~g/dL \cite{friebelDeterminationComplexRefractive2005}. En este trabajo se optó por emplear las propiedades ópticas de la oxihemoglobina, dado que las condiciones hematológicas analizadas se caracterizan principalmente por alteraciones en el tamaño, volumen y concentración de hemoglobina de los eritrocitos, sin implicar necesariamente variaciones en el estado de oxigenación de la hemoglobina~\cite{keohaneRodaksHematologyClinical2015}. Los valores de CHCM considerados fueron seleccionados con el propósito de representar los distintos escenarios patológicos descritos en la sección anterior. La parte real del índice de refracción $\eta_{\text{Hb}}$ para todas las concentraciones se reporta en términos de la relación \cite{friebelModelFunctionCalculate2006}
%
\begin{equation}
	\eta_{\text{Hb}}(\lambda,c_{\text{Hb}})=n_{\text{H}_{2}\text{O}}(\lambda)\,[\beta(\lambda)c_{\text{Hb}}+1],
	\label{eq:IR_eri}
\end{equation}
%
 de $n_{\text{Hb}}$ con $c_{\text{Hb}}$. En cuanto a la parte imaginaria, para la concentración de 28.7 g/dL se dispone de datos en el intervalo espectral de 250 nm a 1100 nm. En contraste, para las concentraciones de 30.6 g/dL y 34 g/dL la información experimental de la parte imaginaria sólo está disponible en el rango de 300 a 600 nm. Para el plasma, se emplearon los datos de la parte real reportados en la Ref. \cite{liuMeasurementRefractiveIndex2019} y de la parte imaginaria en la Ref.~\cite{meinkeOpticalPropertiesPlatelets2007}, ambos en el intervalo espectral comprendido entre 400 nm y 750 nm.
donde $n_{\text{H}_{2}\text{O}}(\lambda)$ es el índice de refracción del agua, $\beta(\lambda)$ es el incremento refractivo específico de la hemoglobina y $c_{\text{Hb}}$ es la concentración de hemoglobina. En la Fig.~\ref{subfig:beta} se muestra la dependencia espectral de $\beta(\lambda)$ en función de la longitud de onda $\lambda$ de la luz incidente, donde los puntos indican los datos experimentales y la línea continua una interpolación. En la región ultravioleta se observa dispersión anómala, mientras que en los intervalos 310–355 nm y 500–1100 nm las variaciones de $\beta$ no son significativas: la desviación estándar del promedio en dichos rangos es del orden de $3.6\times10^{-5}$, comparable con el error experimental reportado para mediciones a una longitud de onda dada \cite{friebelModelFunctionCalculate2006}. En la Fig.~\ref{subfig:refrac} se presentan las curvas de la parte real del índice de refracción obtenidas a partir de la Ec.~\eqref{eq:IR_eri} para las distintas concentraciones: la curva verde para 28.7~g/dL, la azul para 30.6 g/dL y la roja para 34.0 g/dL en función de la longitud de onda $\lambda$ de la luz incidente; a largo de este trabajo se mantendrá el código  de color asociado a la concentración de CHCM. La curva gris punteada representa la parte real del índice de refracción del agua, incluida como referencia. La comparación entre las curvas correspondientes a distintas concentraciones de hemoglobina muestran la dependencia lineal
%
\begin{figure}[t]
	\centering
	\sidesubfloat[First image]{\hspace{-0.5cm}{
			\includegraphics[width=7.5cm]{2-Eritrocitos/beta.pdf}\label{subfig:beta}}}\hspace{0.1cm}
	\sidesubfloat[Second image]{\hspace{-0.4cm}{	\includegraphics[width=7.5cm]{2-Eritrocitos/refractiveIndex.pdf}\label{subfig:refrac}}}\vspace{-0.2cm}
	\caption{Datos empleados reportados por Meinke \textit{et al}. \cite{meinkeOpticalPropertiesPlatelets2007} para el índice de refracción de la hemoglobina a distintas concentraciones como función de la longitud de onda $\lambda$. \textbf{a)}~Incremento refractivo específico $\beta(\lambda)$ de la hemoglobina. Los marcadores corresponden a datos experimentales y la línea continua a una interpolación. \textbf{b)} Parte real del índice de refracción de la hemoglobina $\eta_{\text{Hb}}$ para distintas concentraciones: 28.7 g/dL, 30.6 g/dL y 34.0~g/dL indicados mediante las líneas continuas verde, azul y roja, respectivamente. La línea gris punteada corresponde a la parte real del índice de refracción del agua obtenido de la Ref. \cite{haleOpticalConstantsWater1973}. }
	\label{fig:dataEri}
\end{figure}
%

Con el fin de disponer de una representación continua y analítica de la función dieléctrica en todo el rango del visible, la parte imaginaria de la función dieléctrica de los eritrocitos y del plasma se ajustó mediante un modelo de osciladores de Lorentz\footnote{En este modelo, los electrones del material se describen como osciladores armónicos amortiguados con frecuencias de resonancia características, sometidos a la acción de un campo electromagnético externo \cite{bohrenAbsorptionScatteringLight2008}.}. Esta elección, aunque no es única, resulta particularmente conveniente porque proporciona una expresión analítica que, por construcción, satisface el principio de causalidad \cite{maggioreModernIntroductionClassical2023} y, por lo tanto, es compatible con las relaciones de KK. El ajuste se realizó por mínimos cuadrados entre los datos experimentales y el modelo propuesto, empleando un número finito de osciladores (ver el Apéndice \ref{section:apendix1}). Para identificar las frecuencias de resonancia que corresponden a los osciladores ajustados, en las siguientes figuras se incluye el valor de la energía correspondiente a la posición espectral de cada resonancia. Este valor está determinado por la expresión $\hbar\omega=hc/\lambda$, donde $\hbar=h/2\pi$ con $h$ la constante de Planck y $\lambda$ la longitud de onda de la luz incidente. En la Fig.~\ref{subfig:ajuste_ery} se muestran los datos experimentales correspondientes a la parte imaginaria de la función dieléctrica ($\text{Im}\{\varepsilon\}$) de un eritrocito sano con una CHCM de 30.6 g/dL, extraídos de la Ref. \cite{friebelDeterminationComplexRefractive2005} (puntos de color azul oscuro), junto con el ajuste obtenido mediante osciladores de Lorentz (línea azul claro continua) en función de la longitud de onda $\lambda$ (eje inferior) y la energía $\hbar\omega$ (eje superior) de la luz incidente. De manera análoga, en la Fig.~\ref{subfig:ajuste_plasma} se presentan los datos experimentales de $\text{Im}\{\varepsilon\}$ del plasma (puntos de color violeta oscuro) y su correspondiente ajuste mediante osciladores de Lorentz (línea violeta claro continua). En ambos casos, las líneas grises verticales indican las frecuencias de resonancia asociadas a los osciladores empleados.
%
\begin{figure}[h!]
	\centering
	\sidesubfloat[First image]{\hspace{-0.5cm}{
		\includegraphics[width=7.5cm]{2-Eritrocitos/ajusteLorentz30.pdf}\label{subfig:ajuste_ery}}}\hspace{0.1cm}
	\sidesubfloat[Second image]{\hspace{-0.5cm}{	\includegraphics[width=7cm]{2-Eritrocitos/ajustePlasma.pdf}\label{subfig:ajuste_plasma}}}\vspace{-0.1cm}
	\caption{Parte imaginaria $\text{Im}\{\varepsilon\}$ en función de la longitud de onda $\lambda$ (eje inferior) y la energía $\hbar\omega$ (eje superior) de la luz incidente. \textbf{a)} Eritrocitos con una concentración de 30.6 g/dL. Los datos experimentales se obtuvieron de la Ref.~\cite{friebelDeterminationComplexRefractive2005} y están indicados mediante marcadores de color azul oscuro. La línea azul claro continua representa el ajuste con lorentzianas obtenido.  \textbf{b)} Plasma, cuyos datos experimentales representados por marcadores de color violeta oscuro se extrajeron de la Ref.~\cite{meinkeOpticalPropertiesPlatelets2007}. La línea violeta claro continua representa el ajuste con lorentzianas obtenido. En ambas figuras las líneas grises verticales indican las frecuencias de resonancia asociadas a los osciladores de Lorentz empleados en el ajuste.}
	\label{fig:ajuste}
\end{figure}
%

Dado que los datos experimentales disponibles cubren únicamente ventanas espectrales finitas, los ajustes empleados no consideran transiciones electrónicas en otras regiones por lo que cálculos con KK no estimarían correctamente la parte real calculada a partir de los datos experimentales. Para mitigar este problema se emplearon las relaciones SSKK, las cuales requieren conocimiento del valor de $\text{Re}\{\varepsilon\}$ o $\text{Im}\{\varepsilon\}$ a una frecuencia específica $\omega_1$, lo que permite reducir los errores asociados al truncamiento espectral \cite{lucariniKramersKronigRelationsOptical2005}. En este trabajo se determinó la energía $\hbar\omega_1$ mediante un ajuste minimizando el error absoluto entre la parte real experimental de la función dieléctrica y la obtenida a partir de las relaciones SSKK en el rango espectral de 400 a 600 nm (ver Apéndice \ref{section:apendix2}). Se obtuvieron valores de 2.85 eV, 2.85 eV  y 2.91 eV para las concentraciones de hemoglobina de 28.7~g/dL, 30.6 g/dL y 34.0 g/dL, respectivamente, y de 2.88~eV para el plasma. En la Fig.~\ref{subfig:epsSSKK_eri} se muestra la comparación entre la parte real de la función dieléctrica ($\text{Re}\{\varepsilon\}$) reconstruida mediante SSKK, indicada mediante la línea azul claro continua y los datos experimentales para una CHCM de 30.6~g/dL indicados mediante los puntos azul oscuro, observándose una desviación promedio de 0.014. La línea azul oscuro continua que se observa es únicamente una guía al ojo. De manera análoga, la Fig. \ref{subfig:epsSSKK_Plasma} muestra la parte real de la función dieléctrica obtenida mediante el análisis de SSKK como una línea violeta claro continua donde se observa una desviación promedio de 0.011 respecto a los datos experimentales indicados mediante puntos violeta oscuro (gráfica principal). Asimismo, se muestra una ampliación de la parte real de la función dieléctrica obtenida mediante el análisis de SSKK en el rango de 350 a 650 nm (gráfica interna). La diferencia observada entre lo obtenido mediante SSKK y los datos experimentales, particularmente en el ultravioleta, se debe a que se ha reportado que el índice de refracción de soluciones de hemoglobina y del plasma presenta un incremento pronunciado en el ultravioleta profundo, asociado principalmente a la absorción del agua \cite{bosschaartLiteratureReviewNovel2014}.  Este comportamiento no se reproduce completamente en el presente análisis, ya que los espectros de absorción disponibles para la hemoglobina y el plasma ---y, por ende, la parte imaginaria de la función dieléctrica empleada--- sólo se extienden hasta aproximadamente 250~nm y 300~nm, respectivamente.  A pesar de esta limitación, la aplicación de las relaciones SSKK permite obtener una reconstrucción cuya desviación promedio con respecto a los datos experimentales es casi 1.5 veces menor que la reportada en la Ref. \cite{bosschaartLiteratureReviewNovel2014}, donde se emplea el valor de $\hbar\omega_1$ igual a 1.5 eV [ver Apéndice \ref{section:apendix2}]. En principio, la extensión de los datos experimentales a un rango espectral más amplio \cite{lucariniKramersKronigRelationsOptical2005} o el uso de múltiples frecuencias de anclaje mediante relaciones MSKK\footnote{Las relaciones de Kramers–Kronig generalizadas con sustracciones múltiples (MSKK) permiten introducir varios valores de la función dieléctrica conocidos para mejorar la precisión del análisis \cite{palmerMultiplySubtractiveKramers1998}.} permitiría mejorar la reconstrucción de la función dieléctrica en la región UV.
%
\vspace{-0.1cm}
\begin{figure}[h]
	\centering
		\sidesubfloat[First image]{\hspace{-0.5cm}{
			\includegraphics[width=7.5cm]{2-Eritrocitos/sskk30.pdf}\label{subfig:epsSSKK_eri}}}\hspace{-0.1cm}
	\sidesubfloat[Second image]{\hspace{-0.5cm}{	\includegraphics[width=7.5cm]{2-Eritrocitos/sskkPlasma.pdf}\label{subfig:epsSSKK_Plasma}}}\vspace{-0.2cm}
	\caption{Comparación entre la parte real de la función dieléctrica $\text{Re}\{\varepsilon\}$ reconstruida mediante las relaciones SSKK (línea color claro continua) y los datos experimentales (marcadores color oscuro) en función de la longitud de onda $\lambda$ de la luz incidente para: \textbf{a)} eritrocitos con una CHCM de 30.6 g/dL, a partir de los datos reportados en la Ref. \cite{friebelDeterminationComplexRefractive2005}, donde la línea azul oscuro continua que se observa es una guía al ojo y \textbf{b)} plasma, empleando los datos experimentales de la Ref. \cite{meinkeOpticalPropertiesPlatelets2007} (gráfica principal), donde la línea violeta oscuro continua que se observa es una guía al ojo. Se muestra además una ampliación de la parte real de la función dieléctrica obtenida mediante el análisis de SSKK en el rango de 350 a 650 nm (gráfica interna).}
	\label{fig:epsSSKK}
\end{figure}
%

Cabe destacar que, aunque para la concentración de 28.7 g/dL se dispone de datos experimentales tanto de la parte real como de la imaginaria en un intervalo espectral relativamente amplio, en este trabajo se empleó de manera consistente la función dieléctrica reconstruida mediante el mismo procedimiento (ajuste con osciladores de Lorentz seguido de SSKK) para las tres concentraciones consideradas. Esta elección garantiza que todas las funciones dieléctricas estén sujetas a los mismos efectos de truncamiento espectral, de modo que las comparaciones posteriores de las propiedades ópticas calculadas reflejen exclusivamente la variación de la CHCM y no diferencias metodológicas en el tratamiento de los datos.

En la Fig.~\ref{fig:epsEri} se muestran $\text{Re}\{\varepsilon\}$ [Fig. \ref{subfig:epsReEri}] e $\text{Im}\{\varepsilon\}$ [Fig. \ref{subfig:epsImEri}] en función de la longitud de onda $\lambda$ de la luz incidente para las distintas concentraciones de hemoglobina consideradas: verde para 28.7~g/dL, azul para 30.6 g/dL y rojo para 34.0 g/dL. En todas las curvas destaca la banda de Soret, característica de la hemoglobina, con un máximo situado típicamente en el intervalo 410–430~nm~\cite{bosschaartLiteratureReviewNovel2014}, cuya intensidad es considerablemente mayor que la de las transiciones en el visible. En esta última región se observan las denominadas bandas Q, correspondientes a la banda $\alpha$ (a mayor longitud de onda) y la banda $\beta$ (a menor longitud de onda), cuyos máximos se localizan de manera general en los rangos aproximados 600–650 nm y 520–580 nm \cite{dayerBandAssignmentHemoglobin2010}, respectivamente, dependiendo de la concentración y del estado de oxigenación\footnote{Estas bandas se originan en transiciones electrónicas $\pi \rightarrow \pi^*$ del anillo de porfirina. Tanto la banda de Soret (UV–Vis) como las bandas Q (visible) son sensibles al estado de oxidación del hierro y a cambios conformacionales de la hemoglobina, por lo que caracterizan su estructura y función \cite{dayerBandAssignmentHemoglobin2010}.}. Las regiones sombreadas en la Fig. \ref{fig:epsEri} señalan las bandas espectrales mencionadas anteriormente. De manera adicional, tanto en la parte real como en la imaginaria se aprecia un relieve alrededor de 450 nm, que coincide con el flanco de la banda de Soret y marca la transición entre el régimen de dispersión anómala en el UV-Vis y el comportamiento más suave característico del visible.
%
\begin{figure}[H]
	\centering
	\sidesubfloat[First image]{\hspace{-0.8cm}{
			\includegraphics[width=0.47\textwidth]{2-Eritrocitos/parteReEri.pdf}\label{subfig:epsReEri}}}\hspace{0.1cm}
	\sidesubfloat[Second image]{\hspace{-0.4cm}{	\includegraphics[width=0.47\textwidth]{2-Eritrocitos/parteImEri.pdf}\label{subfig:epsImEri}}}\vspace{-0.2cm}
	\caption{Función dieléctrica de eritrocitos en función de la longitud de onda $\lambda$ de la luz incidente para distintas concentraciones de hemoglobina, indicadas mediante líneas continuas verde, azul y roja correspondientes a 28.7 g/dL, 30.6~g/dL y 34 g/dL, respectivamente. Las regiones sombreadas señalan las bandas espectrales que dominan la respuesta óptica de los eritrocitos en el UV-Vis. \textbf{a)} Parte real de la función dieléctrica. \textbf{b)} Parte imaginaria de la función dieléctrica.}
	\label{fig:epsEri}
\end{figure}
%
